% SOURCE docs/release/2026-10-29/source/docs/writing-repair/2026-10-29/src/ce2960f99921f249.tex CHAR 0
\chapter{Continuous time, observation events and structural dynamics}
\label{r3:continuous}
A credit line can change on a decision date, a payment can arrive at any
instant, and a financial statement can be entered weeks after the period
it describes. Putting all three on a monthly grid is a modeling choice.
The event clock in Chapter~\ref{r3:longitudinal} must extend to continuous
time without confusing a changing customer state, the times at which the
bank observes it, and the times at which the bank acts.

\section{A process equation acquires causality through modular interventions}
\label{r3:dynamic-semantics}
On a finite horizon $[0,T]$, consider the structural diffusion
\begin{equation}
 dX_t=b(X_t,A_t,t)\,dt+\Sigma(X_t,A_t,t)\,dW_t,
 \qquad X_0\sim P_0.
 \label{r3:structural-diffusion}
\end{equation}
Here $X_t$ is the scientific state, $A_t$ a predictable action, and $W_t$
an exogenous Brownian motion independent of the initial state. A specified
joint noise law, nonanticipating coefficients, and existence and uniqueness
under the proposed actions are part of the model. Global Lipschitz and
linear-growth conditions supply a familiar sufficient well-posedness route.
The causal assertion is that replacing the mechanism for $A$ leaves the
other structural mechanisms and exogenous law invariant. A fitted
conditional drift without that assertion describes observational dynamics.

A deterministic intervention can specify the whole action path $a(t)$.
A feedback intervention specifies a predictable rule using information
available immediately before $t$. A jump in an account limit, a continuing
rate incentive, and a continuously clamped balance are different
interventions even if their observed values agree at one instant.
Intervening on an accounting total may require an operational action that
changes several underlying flows. The mathematical substitution must have
a feasible banking interpretation.

A dynamic SCM treats each coordinate's entire path as a random element.
This allows feedback between processes while excluding dependence on their
future realizations. \citet[Definitions 4--7 and Sections 4--5]{boeken2024}
develops this path-space connection and defines time splitting and
subsampling. The later \citet[Definitions 3--5 and Theorems 3, 5--6]{sde2026}
distinguishes essential unique solvability for every subsystem, preserving
it under marginalization. These conditions give well-defined causal
laws; they are not inferred from a low prediction error.

The distinction between a graphical Markov property and intervention
calculus matters here. The 2026 source proves a $\sigma$-separation Markov
property under its general solvability conditions and a stronger
$d$-separation property for its specified additive-noise class. Its
Theorem~12 establishes the $\sigma$-separation do-calculus, with
absolute-continuity conditions for the relevant conditional laws.
It explicitly does \emph{not} deduce the cyclic $d$-separation do-calculus
from the stronger Markov property; that extension is conjectural in the
archived text. An implementation may therefore not replace every path-space
separation test by an ordinary DAG routine. On acyclic graphs the two
separations coincide.

\section{Exact sampling and an observational alias}
\label{r3:diffusion-alias}
For a linear structural system with a fixed action $a_k$ between consecutive
sampling times,
\[
 dX_t=(FX_t+Ba_k)\,dt+L\,dW_t,\qquad t_k<t\le t_{k+1},
\]
write $\Delta_k=t_{k+1}-t_k$. Multiplying by $e^{-Ft}$ and integrating gives
\begin{align}
 X_{k+1}&=\Phi_{\Delta_k}X_k+D_{\Delta_k}a_k+\eta_k,\\
 \Phi_\Delta&=e^{F\Delta},\qquad
 D_\Delta=\int_0^\Delta e^{Fs}B\,ds,\qquad
 Q_\Delta=\Var(\eta_k)=\int_0^\Delta e^{Fs}LL'e^{F's}\,ds .
 \label{r3:exact-sampling}
\end{align}
The stochastic integral is Gaussian with mean zero; independent Brownian
increments make it independent of the past. An exogenous sampling schedule
therefore yields an exact Gaussian state transition. Euler replacement
$\Phi_\Delta=I+F\Delta$ is an approximation with a step-size error.
Using the same transition for irregularly spaced observations does not
represent the same continuous process.

Even exact discrete transitions need not identify $F$. Consider
\[
 F_\omega=-aI+\omega J,\qquad
 J=\begin{pmatrix}0&-1\\1&0\end{pmatrix},\qquad L=\sigma I,\quad a>0 .
\]
Since $J^2=-I$,
\[
 \Phi_\Delta=e^{-a\Delta}
 \{\cos(\omega\Delta)I+\sin(\omega\Delta)J\},\qquad
 Q_\Delta=\frac{\sigma^2}{2a}(1-e^{-2a\Delta})I.
\]
With no historical control input, $\omega$ and
$\omega+2\pi/\Delta$ give identical transition laws on the regular grid,
for any common initial distribution. Their drift and off-grid responses
differ. More observations on that same grid cannot resolve this alias.

The difference also affects a new sustained action. Set
$a=\sigma=\Delta=1$ and suppose, as extra structural knowledge, that
$B=(1,0)'$ is known. Starting from zero, a unit action held throughout the
interval gives
\[
 D_1(0)=\begin{pmatrix}1-e^{-1}\\0\end{pmatrix},\qquad
 D_1(2\pi)=\frac{1-e^{-1}}{1+4\pi^2}
             \begin{pmatrix}1\\2\pi\end{pmatrix}.
\]
The first-coordinate responses are approximately $.6321$ and $.0156$.
These project constructions are checked against independent
matrix-exponential and integration calculations. They isolate drift
ambiguity even with the control direction supplied; the historical
zero-input law also lacks direct evidence for the new action.
Varied sampling intervals or informative exogenous controls can distinguish
these particular models, but that requires a new identification argument.

Subsampling marginalizes unobserved path segments. It can create
confounding paths, eliminate short-lived responses and destroy conditional
independences among complete paths. This is the operation in
\citet[Definition 9 and Example 10]{sde2026}, not a neutral change of file
format. Independent innovations in a structural diffusion need not yield
independent contemporaneous residual coordinates after integrating out
intermediate dynamics.

\section{Equilibrium and an intervention on a feedback rule}
\label{r3:equilibrium}
The steady balance response to a bank's pricing rule differs from its
response to a clamped offered rate. For a small illustrative model, let
$x$ denote scaled deposits and $r$ a rate deviation:
\[
 \dot x=-x+br+u,\qquad
 \dot r=-\kappa(r-r_0+cx),\qquad b,c,\kappa>0.
\]
The customer disturbance $u$ is constant for a trajectory. The bank lowers
its incentive as deposits increase; $r_0$ is the intercept of that rule.
The system matrix has negative trace $-(1+\kappa)$ and positive determinant
$\kappa(1+bc)$, so it converges to its unique equilibrium:
\[
 x^*=\frac{br_0+u}{1+bc},\qquad r^*=r_0-cx^*.
\]
Changing the rule intercept has derivative
$\partial x^*/\partial r_0=b/(1+bc)$.
Fixing $r(t)=\bar r$ removes its feedback equation and instead yields
$x^*=b\bar r+u$, with derivative $b$.
If only $u$ varies across observed equilibria while $r_0$ is fixed,
the observational slope of $x^*$ on $r^*$ is $-1/c$.
A negative fitted rate--balance relation is compatible with a positive
structural rate response in this specified model.

This calculation requires convergence under the relevant intervention.
Setting derivatives to zero gives candidate equilibrium equations; it
does not establish that actual paths reach them.
\citet[Definition 4.9, Theorems 4.11 and 4.18]{bongers2022}
formalizes equilibration and its commutation with steady interventions
under the stated convergence conditions. Nonconverging paths, multiple
solutions, and selection of only equilibrated units need separate analysis.
A stationary diffusion generally fluctuates forever: an invariant
distribution is not a pathwise equilibrium obtained by setting Brownian
forcing to zero.

A bank action can also change competitors' responses and market clearing.
The example exposes the feedback distinction but is not a model of an
entire banking market. A portfolio calculation holding competitor rates
fixed and a market equilibrium counterfactual must be different product
requests, with different maintained mechanisms.

\section{An event intensity is a conditional law for a clock}
\label{r3:event-intensity}
Let $N^k_t$ count events of type $k$, with no simultaneous events,
no explosions on $[0,T]$, and predictable intensities $\lambda^k_t$
relative to a declared history $\mathcal F_{t-}$.
Its compensator is $\Lambda^k_t=\int_0^t\lambda^k_s\,ds$, and
$M^k_t=N^k_t-\Lambda^k_t$ is a martingale under suitable integrability.
Examples include payments, delinquency transitions, contact attempts,
report arrivals and closure. A marked process also records the amount,
action or covariate revealed at an event.

Conditional on the history, the probability of one type-$k$ event in a
short interval is $\lambda^k_t\,dt+o(dt)$; the probability of no event is
$1-\sum_k\lambda^k_t\,dt+o(dt)$. Multiplying along a path and taking a
limit gives the intensity-dependent likelihood
\[
 \mathcal L(\lambda)\ \propto\
 \exp\left\{\sum_k\int_0^T\log\lambda^k_t\,dN^k_t
                   -\sum_k\int_0^T\lambda^k_t\,dt\right\}.
\]
The integrated hazard accounts for every interval without an event.
A likelihood fitted only to recorded event rows discards that information.

Suppose an intervention replaces treatment-event intensities, indexed
by $\mathcal A$, by externally specified predictable $\lambda^{k,g}_t$.
Assume causal invariance of the other local laws. Cancelling their
likelihood factors gives
\begin{equation}
 W_T^g=
 \exp\left\{\sum_{k\in\mathcal A}\left[
 \int_0^T\log\frac{\lambda^{k,g}_t}{\lambda^k_t}\,dN^k_t
 -\int_0^T(\lambda^{k,g}_t-\lambda^k_t)\,dt\right]\right\},
 \qquad \E^g[Y]=\E[W_T^gY].
 \label{r3:event-weight}
\end{equation}
For zero target intensity, use the product form: a prohibited event makes
the weight zero. The equality requires an actual likelihood ratio of
mean one, support and integrability, as well as causal validity.
Bounded intensities and bounded intensity ratios on a finite horizon
provide a useful sufficient regular regime. A nonnegative local
martingale alone can lose expectation and need not define the intended
probability law. This reconstructs
\citet[Definition 2 and Proposition 1]{roysland2024}.

The event-ratio product alone is insufficient. For a Poisson treatment
clock of rate $\lambda=1$ changed to $\mu=2$ over unit time, the correct
weight is $e^{-1}2^{N_1}$. Its expectation is one, and
$\E[e^{-1}2^{N_1}N_1]=2$. Omitting $e^{-1}$ gives mean weight $e$.
Normalizing weights in this constant-rate example happens to cancel the
error; with state- or history-dependent rates the missing integral varies
across trajectories and that cancellation fails.

The filtration is a causal restriction. If a collector contacts a customer
using a private distress signal $S_t$, the observable intensity
$\E[\lambda_t(S_t,\mathcal F_{t-})\mid\mathcal O_{t-}]$ describes recorded
event times but is not automatically a valid causal replacement mechanism.
Marginalizing $S$ may change which other intensities remain invariant.
The graphical eliminability conditions in
\citet[Theorem 1 and Section 6]{roysland2024} address this issue.
Suppressing censoring also needs a defensible observation-only intervention;
independent censoring as a probabilistic statement alone does not establish it.

\section{Assignments at visits, visit schedules and deterministic regimes}
\label{r3:event-regimes}
Suppose $N^v$ counts decision opportunities. At each event the bank assigns
a mark $A$ with probability $\pi(a\mid\mathcal F_{t-})$.
Changing the action distribution at existing opportunities contributes
$\prod_{t:\Delta N^v_t=1}g(A_t\mid\mathcal F_{t-})/
\pi(A_t\mid\mathcal F_{t-})$ to the weight. Changing the opportunity
intensity requires the additional waiting-time and event terms in
Equation~\eqref{r3:event-weight}. Contacting every existing visitor
differently and doubling contact frequency are different experiments.

\citet[Definitions 1--2]{rytgaard2021} explicitly distinguishes these
targets. Its main efficient-gradient derivation changes treatment marks
and censoring while retaining the treatment-monitoring clock among
non-intervened mechanisms. That clock contributes a nuisance-score term.
Reusing the score for an intervention that fixes the clock is incorrect.
A deterministic action mark at a supported visit can be identified even
though forcing a new visit at an exact calendar instant is singular
relative to a clock with a continuous event-time density.

A different route starts from a deterministic, nonanticipating regime on
the event path. In the 2026 preprint
\citet[Definition 3 and Theorem 1]{ryalen2026}, let $\tau^a$ be the first
deviation of the observed treatment path from the regime. Write
$D_t=1\{\tau^a\le t\}$ and let $\Lambda^D$ be its compensator.
The source constructs
\[
 K_t=-\int_0^t\frac{d(D_s-\Lambda^D_s)}{1-\Delta\Lambda^D_s},
 \qquad dW_t=W_{t-}\,dK_t,\quad W_0=1.
\]
Its positivity condition controls
$\int_0^T d\Lambda^D_s/(1-\Delta\Lambda^D_s)$.
Consistency holds up to deviation. Exchangeability requires the same
deviation compensator after enlarging the history by the future potential
outcome path; it is a martingale restriction, not a symmetric conditional
independence between two time-indexed variables. Mean-one likelihood-ratio
regularity is an additional requirement.

The proof is short once these premises hold. $W_t$ is zero on deviated
paths, so consistency gives $\E[W_tY_t]=\E[W_tY_t^a]$.
The weight remains a martingale in the enlarged filtration, yielding
$\E[W_tY_t^a]=\E[W_0Y_t^a]=\E[Y_t^a]$ for bounded outcomes.
On a fixed discrete decision schedule, these adherence weights reduce to
products of inverse action probabilities.

This does not make an arbitrary deterministic event path supported by a
continuous clock. The source distinguishes its regime-specific g-formula
distributions from the larger class of measures obtained by changing
compensators while preserving other local laws
\citep[Definition 5, Remark 2 and Theorem 2]{ryalen2026}.
The routes encode different intervention classes and premises. A product
must state which one justifies a request, rather than attach
\emph{continuous-time g-formula} to any hazard model.

\section{A martingale gradient shows what an estimator needs}
\label{r3:event-gradient}
For a transparent restricted construction, take a finite, fully observed
marked jump system with bounded intensities and a bounded terminal payoff.
Let $g$ replace a specified subset of event intensities; the remaining
mechanisms are indexed by $\mathcal E$. Write $H_{t-}$ for the full
event history, and $V^g(t,h)$ for conditional expected terminal payoff
under the target law from that history. Let $\Delta_kV^g(t,h)$ be the
change in continuation value on appending environmental event $k$ at
time $t$. There are no additional continuously observed noises or random
marks omitted from this event alphabet.

In the unrestricted model for these history-dependent intensities,
perturbing environmental intensity $k$ with log-score $s_k$ gives likelihood
score $\int s_k\,dM^k$. Under the target law, conditioning the terminal
payoff converts its covariance with that score to the expected integral
of $\Delta_kV^g s_k\lambda^k\,dt$. Changing measure back to the observed
history introduces $W^g_{t-}$. Distinct event martingales are orthogonal
because there are no simultaneous events. Thus the project gradient is
\begin{equation}
 \phi_g=V^g(0,H_0)-\psi_g+
 \sum_{k\in\mathcal E}\int_0^T W^g_{t-}\Delta_kV^g(t,H_{t-})
       \{dN^k_t-\lambda^k_t\,dt\}.
 \label{r3:jump-gradient}
\end{equation}
Treatment-intensity scores have zero derivative because $g$ is fixed.
Baseline scores recover the first term. Square integrability and the
specified tangent-space decomposition justify the representation.
A Markov restriction reduces the tangent space; this formula should not
be labeled efficient for that smaller model without its own projection.

For a single absorbing event of constant hazard $\lambda$, take
$Y=1\{N_T=0\}$. Before the event,
$V(t)=e^{-\lambda(T-t)}$ and $\Delta V(t)=-e^{-\lambda(T-t)}$.
Direct integration of $\int\Delta V\,dM$ gives $Y-e^{-\lambda T}$:
on an event at $u<T$ it is
\[
 -e^{-\lambda(T-u)}
 +\int_0^u\lambda e^{-\lambda(T-t)}dt=-e^{-\lambda T};
\]
on no event it is $1-e^{-\lambda T}$.
This pathwise check connects the martingale expression to the ordinary
gradient for a mean.

The general covariate-mark model needs both mark innovations and
monitoring, treatment-opportunity and outcome martingales.
\citet[Theorem 1 and equations (9)--(13)]{rytgaard2021} derives that
richer gradient and a second-order remainder. Its Theorem~2 requires
an empirical score of $o_p(n^{-1/2})$, a remainder of that order, and
a Donsker plus mean-square convergence condition for the estimated score.
Solving the targeting equation alone is insufficient. Cross-fitting
requires a different empirical-process argument; the source theorem does
not validate an arbitrary cross-fitted neural intensity fit.

Here $n$ counts independent customer trajectories, not transaction rows.
Common calendar shocks create another dependence level. Numerical
integration, event-time precision, intensity estimation and Monte Carlo
error need separate tolerances relative to financial-target uncertainty.

\section{An absence of observations is evidence about hidden state}
\label{r3:no-event-filter}
Let a finite hidden state have row generator $Q$ and observation-event
rate $\lambda_s$ in state $s$. Conditional on an event, the mark has
density $f_s(y)$. Assume event generation does not change the state and
that state jumps and observation events do not coincide.
For a row vector $\alpha$ of unnormalized state probabilities, a short
interval without an event gives
\[
 \alpha(t+dt)=\alpha(t)\{I+Q\,dt-\operatorname{diag}(\lambda)\,dt\}
               +o(dt).
\]
For a constant generator and rates over an interval, it follows that
\begin{align}
 \alpha(t+\Delta)&=\alpha(t)
       e^{\{Q-\operatorname{diag}(\lambda)\}\Delta}
                    &&\text{if no event occurs},\\
 \alpha(t^+)&=\alpha(t^-)
       \operatorname{diag}\{\lambda_s f_s(y)\}
                    &&\text{at an event with mark }y .
 \label{r3:event-filter}
\end{align}
The sum of the final probabilities is the likelihood density of the
observed schedule and marks. Dividing by that sum gives the filter.
Normalizing the differential equation yields
\[
 \dot p_s=(pQ)_s-p_s(\lambda_s-\bar\lambda),
 \qquad \bar\lambda=\sum_s p_s\lambda_s .
\]
The second term shifts probability away from states in which an event
would have been more likely, even while recorded covariates stay unchanged.

For a static two-state customer, let prior distress probability be $1/2$,
with reporting rates $\lambda_0=.1$ and $\lambda_1=1$ per month.
After a month with no report,
\[
 P(S=1\mid N_1=0)=\frac{e^{-1}}{e^{-.1}+e^{-1}}
                 =\frac{1}{1+e^{.9}}\approx .2891.
\]
An unmarked report immediately afterward updates this to approximately
$.8026$. Ignoring the preceding silence gives $1/(1+.1)\approx .9091$.
Both use the same report and initial prior; their information sets differ.

A missingness flag can represent silence but does not estimate its
state-dependent intensity. If scheduling depends only on recorded history,
its factor can sometimes be conditioned away. If it depends on private
state, it belongs in the joint likelihood. Under an intervention making
visits externally scheduled, the historical state-dependent scheduling
factor must be replaced. Retaining it in the target filter manufactures
information from an intervention that no longer carries that information.

The computational contract includes continuous dynamics, action and
observation clocks, event marks, time units, the sampling operator and
information available to each policy. An ODE/SDE solver, event likelihood
and causal identification procedure solve different parts of that contract.
Their numerical success must remain separate from evidence for its
structural assumptions.

